\documentclass[10pt,a4paper]{amsart}
\usepackage{amsmath,amssymb,amsthm,graphicx}
\usepackage{verbatim}
\newenvironment{reference}{\comment}{\endcomment}
%\newenvironment{reference}{}{}
\newenvironment{slogan}{\comment}{\endcomment}
\newenvironment{history}{\comment}{\endcomment}
% For commutative diagrams we use Xy-pic
\usepackage[all]{xy}

% We use 2cell for 2-commutative diagrams.
\xyoption{2cell}
% Theorem environments.
%
\theoremstyle{plain}
\newtheorem{theorem}[subsection]{Theorem}
\newtheorem{proposition}[subsection]{Proposition}
\newtheorem{lemma}[subsection]{Lemma}

\theoremstyle{definition}
\newtheorem{definition}[subsection]{Definition}
\newtheorem{example}[subsection]{Example}
\newtheorem{exercise}[subsection]{Exercise}
\newtheorem{situation}[subsection]{Situation}

\theoremstyle{remark}
\newtheorem{remark}[subsection]{Remark}
\newtheorem{remarks}[subsection]{Remarks}

\numberwithin{equation}{subsection}

% Macros
%
\def\lim{\mathop{\mathrm{lim}}\nolimits}
\def\colim{\mathop{\mathrm{colim}}\nolimits}
\def\Spec{\mathop{\mathrm{Spec}}}
\def\Hom{\mathop{\mathrm{Hom}}\nolimits}
\def\Ext{\mathop{\mathrm{Ext}}\nolimits}
\def\SheafHom{\mathop{\mathcal{H}\!\mathit{om}}\nolimits}
\def\SheafExt{\mathop{\mathcal{E}\!\mathit{xt}}\nolimits}
\def\Sch{\mathit{Sch}}
\def\Mor{\mathop{\mathrm{Mor}}\nolimits}
\def\Ob{\mathop{\mathrm{Ob}}\nolimits}
\def\Sh{\mathop{\mathit{Sh}}\nolimits}
\def\NL{\mathop{N\!L}\nolimits}
\def\CH{\mathop{\mathrm{CH}}\nolimits}
\def\proetale{{pro\text{-}\acute{e}tale}}
\def\etale{{\acute{e}tale}}
\def\QCoh{\mathit{QCoh}}
\def\Ker{\mathop{\mathrm{Ker}}}
\def\Im{\mathop{\mathrm{Im}}}
\def\Coker{\mathop{\mathrm{Coker}}}
\def\Coim{\mathop{\mathrm{Coim}}}

% Boxtimes
%
\DeclareMathSymbol{\boxtimes}{\mathbin}{AMSa}{"02}

%
% Macros for moduli stacks/spaces
%
\def\QCohstack{\mathcal{QC}\!\mathit{oh}}
\def\Cohstack{\mathcal{C}\!\mathit{oh}}
\def\Spacesstack{\mathcal{S}\!\mathit{paces}}
\def\Quotfunctor{\mathrm{Quot}}
\def\Hilbfunctor{\mathrm{Hilb}}
\def\Curvesstack{\mathcal{C}\!\mathit{urves}}
\def\Polarizedstack{\mathcal{P}\!\mathit{olarized}}
\def\Complexesstack{\mathcal{C}\!\mathit{omplexes}}
% \Pic is the operator that assigns to X its picard group, usage \Pic(X)
% \Picardstack_{X/B} denotes the Picard stack of X over B
% \Picardfunctor_{X/B} denotes the Picard functor of X over B
\def\Pic{\mathop{\mathrm{Pic}}\nolimits}
\def\Picardstack{\mathcal{P}\!\mathit{ic}}
\def\Picardfunctor{\mathrm{Pic}}
\def\Deformationcategory{\mathcal{D}\!\mathit{ef}}
\usepackage{cleveref}
\crefname{originallemma}{original lemma}{original lemmas}
\Crefname{originallemma}{Original lemma}{Original lemmas}
\makeatletter
\newlabel{lemma-flat-bc-surjection}{{lemma-flat-bc-surjection}{source}{}{sourcepointer-1}{}}
\newlabel{lemma-flat-bc-surjection@cref}{{[originallemma][0][]lemma-flat-bc-surjection}{[1][0][]source}{}{}{}}
\makeatother
\crefname{originalreference}{original reference}{original references}
\Crefname{originalreference}{Original reference}{Original references}
\makeatletter
\newlabel{more-algebra-lemma-double-base-change}{{more-algebra-lemma-double-base-change}{source}{}{sourcepointer-2}{}}
\newlabel{more-algebra-lemma-double-base-change@cref}{{[originalreference][0][]more-algebra-lemma-double-base-change}{[1][0][]source}{}{}{}}
\makeatother
\crefname{originalsection}{original section}{original sections}
\Crefname{originalsection}{Original section}{Original sections}
\makeatletter
\newlabel{section-trivial}{{section-trivial}{source}{}{sourcepointer-3}{}}
\newlabel{section-trivial@cref}{{[originalsection][0][]section-trivial}{[1][0][]source}{}{}{}}
\makeatother
\title{The upper shriek functors compose up to a canonical isomorphism on bounded below complexes}
\author{The Stacks Project Authors}
\date{stacks-project:dualizing.tex:local-snapshot-20261012; GNU FDL 1.2 or later}
\begin{document}
\maketitle
\noindent\textit{Editorial note.} All delivered mathematical text is the author's own, copied from the pinned Stacks Project chapter source sources/stacks-dualizing/pinned-original-source.tex; no statement or proof step was rewritten, reordered or removed.\par\medskip
\noindent\textit{Editorial note.} Preamble provenance: the Stacks Project keeps one single shared preamble for every chapter; this card takes it directly from the staged sources/stacks-dualizing/preamble.tex as three verbatim spans (lines 8-12, 14-18 and 154-217), so no derived file was written and nothing was added or redefined. The wrapper class and the base packages are supplied by the assembler.\par\medskip
\noindent\textit{Editorial note.} Bibliography provenance: neither the counted statement nor its proof contains a citation (this lane rescans the whole span before writing the card), so the card carries no bibliography and no entry was reconstructed.\par\medskip
\noindent\textit{Editorial note.} Selection provenance: this card was opened under the file-based control contract (CONTROL.md), with the mechanical pre-card gate run on its statement and proof interval before the card was written.\par\medskip
\noindent\textit{Allowed standard prerequisites (editorial).}\begin{itemize}
\item The base change isomorphism for the upper shriek of a surjection of finite type algebras over a Noetherian ring. Invoked inside the counted proof (source lines 4794-4796). Identifies the shriek of the polynomial extension with the tensor of the shriek after base change in the bounded below case.
\item The identification of the upper shriek of a composition with the composite of the upper shrieks. Invoked inside the counted proof (source lines 4797-4799). Identifies the two ways of writing the shriek of the composite surjection through adjointness to restriction.
\item The double base change isomorphism for the shriek of finite type homomorphisms. Invoked at the end of the counted proof (source lines 4811-4812). Identifies the base changed shriek term with the shriek computed over the source.
\end{itemize}\medskip
\ifdefined\hypertarget\hypertarget{sourcepointer-1}{}\fi
\noindent\textit{Source reference.} Original lemma lemma-flat-bc-surjection: sources/stacks-dualizing/pinned-original-source.tex, source lines 2726--2726. Navigation pointer only: the counted text refers to the Stacks item lemma-flat-bc-surjection, which is an existing labeled result of this same chapter and is not part of this card. The pointer carries the target's own label as its number and its definition line as its span; its content is not reproduced here.\par\medskip
\ifdefined\hypertarget\hypertarget{sourcepointer-2}{}\fi
\noindent\textit{Source reference.} Original reference more-algebra-lemma-double-base-change: sources/stacks-dualizing/pinned-original-source.tex, source lines 4695--4695. Navigation pointer only: the counted text refers to the Stacks item more-algebra-lemma-double-base-change, which belongs to another chapter of the project and is not part of this card. The pointer carries the target's own label as its number because this chapter's staged snapshot does not contain the other chapter's tag table; its content is not reproduced here.\par\medskip
\ifdefined\hypertarget\hypertarget{sourcepointer-3}{}\fi
\noindent\textit{Source reference.} Original section section-trivial: sources/stacks-dualizing/pinned-original-source.tex, source lines 2416--2416. Navigation pointer only: the counted text refers to the Stacks item section-trivial, which is an existing labeled result of this same chapter and is not part of this card. The pointer carries the target's own label as its number and its definition line as its span; its content is not reproduced here.\par\medskip
\begin{lemma}
\label{lemma-composition-shriek-algebraic}
Let $A \xrightarrow{a} B \xrightarrow{b} C$ be finite type homomorphisms of
Noetherian rings. Then there is a transformation of functors
$b^! \circ a^! \to (b \circ a)^!$ which is an isomorphism on $D^+(A)$.
\end{lemma}

\begin{proof}
Choose a polynomial ring $P = A[x_1, \ldots, x_n]$ over $A$
and a surjection $P \to B$. Choose elements $c_1, \ldots, c_m \in C$
generating $C$ over $B$. Set $Q = P[y_1, \ldots, y_m]$ and
denote $Q' = Q \otimes_P B = B[y_1, \ldots, y_m]$.
Let $\chi : Q' \to C$ be the surjection sending $y_j$ to $c_j$.
Picture
$$
\xymatrix{
& Q \ar[r]_{\psi'} & Q' \ar[r]_\chi & C \\
A \ar[r] & P \ar[r]^\psi \ar[u] & B \ar[u]
}
$$
By Lemma \ref{lemma-flat-bc-surjection} for $M \in D(P)$ we have an arrow
$\psi^!(M) \otimes_B^\mathbf{L} Q' \to (\psi')^!(M \otimes_P^\mathbf{L} Q)$
which is an isomorphism whenever $M$ is bounded below. Also
we have $\chi^! \circ (\psi')^! = (\chi \circ \psi')^!$ as both
functors are adjoint to the restriction functor $D(C) \to D(Q)$
by Section \ref{section-trivial}. Then we see
\begin{align*}
b^!(a^!(K))
& =
\chi^!(\psi^!(K \otimes_A^\mathbf{L} P)[n] \otimes_B^\mathbf{L} Q)[m] \\
& \to
\chi^!((\psi')^!(K \otimes_A^\mathbf{L} P \otimes_P^\mathbf{L} Q))[n + m] \\
& =
(\chi \circ \psi')^!(K\otimes_A^\mathbf{L} Q)[n + m] \\
& =
(b \circ a)^!(K)
\end{align*}
where we have used in addition to the above
More on Algebra, Lemma \ref{more-algebra-lemma-double-base-change}.
\end{proof}

\end{document}
